% GENERATED FILE. Edit canonical lessons, then run npm run generate:books.
% canonical-source-hash: fnv1a64:173ad8e0c634408c
% canonical-lessons: GU-C219-gajvu, GU-C219-camakvu, GU-C219-chavavu, GU-C219-kara, GU-C219-mavthu

\chapter{Thunder and hail}
\label{ch:gu-thunder-and-hail}
\hlchaptermodality{219}

\begin{chapteropening}
\textbf{By the end of this chapter:} \emph{I can say that it will thunder and that hail will fall.}

Complete the last lesson of chapter 219: I can say that it will thunder and that hail will fall.
\end{chapteropening}

\section[\texorpdfstring{\gu{ગાજવું} (gājvũ)}{ગાજવું (gājvũ)}]{\gu{ગાજવું} (gājvũ) — to thunder, to rumble}

\label{lesson:GU-C219-gajvu}



\begin{quote}
Before the new one: say the Gujarati for sunset, then the Gujarati for daylight.
\end{quote}

\subsection*{You'll want to know: \gu{ગાજવું}}
\textbf{\gu{ગાજવું}} — \emph{gājvũ} — \textquotedblleft{}to thunder, to rumble\textquotedblright{}.

\begin{quote}
\textbf{\gu{રાત્રે} \gu{વાદળ} \gu{ગાજશે}.} — \emph{rātre vādaḷ gājśe} — \textquotedblleft{}The clouds will thunder at night.\textquotedblright{}
\end{quote}

\begin{cousinweb}
\textbf{\gu{ગાજવું}} (\emph{gājvũ}) is a tadbhava, a word inherited from Sanskrit and reshaped by sound change. It comes from \emph{garjati}, \textquotedblleft{}thunders, roars\textquotedblright{} (CDIAL 4046).
\end{cousinweb}

\begin{grammarlens}[title={What you've built}]
One more word for everyday life at A2.
\end{grammarlens}

\subsection*{Your turn}
\begin{itemize}
\raggedright
  \item \emph{Say it:} \emph{gājvũ}
  \item \emph{Say it:} \emph{gājvũ}, once more
  \item \emph{Say it:} \emph{rātre vādaḷ gājśe} — The clouds will thunder at night.
  \item \emph{Recall:} say the Gujarati for sunset, then the Gujarati for daylight, then say \emph{gājvũ} again
\end{itemize}

\begin{tcolorbox}[breakable,colback=teal!4,colframe=teal!35!black,title={Before you move on}]
What does \gu{ગાજવું} mean? (\textquotedblleft{}To thunder, to rumble\textquotedblright{}.) Say it once more.
\end{tcolorbox}
\section[\texorpdfstring{\gu{ચમકવું} (camakvũ)}{ચમકવું (camakvũ)}]{\gu{ચમકવું} (camakvũ) — to shine, to flash}

\label{lesson:GU-C219-camakvu}



\begin{quote}
Before the new one: say the Gujarati for daylight, then the Gujarati for to thunder.
\end{quote}

\subsection*{You'll want to know: \gu{ચમકવું}}
\textbf{\gu{ચમકવું}} — \emph{camakvũ} — \textquotedblleft{}to shine, to flash\textquotedblright{}.

\begin{quote}
\textbf{\gu{આકાશમાં} \gu{વીજળી} \gu{ચમકશે}.} — \emph{ākāśmā̃ vījḷī camakśe} — \textquotedblleft{}Lightning will flash in the sky.\textquotedblright{}
\end{quote}

\begin{cousinweb}
\textbf{\gu{ચમકવું}} (\emph{camakvũ}) comes from \emph{*camakka}, \textquotedblleft{}sudden movement\textquotedblright{}, a reconstructed form; CDIAL glosses the Gujarati verb \textquotedblleft{}to glitter, be startled\textquotedblright{} (CDIAL 4676).
\end{cousinweb}

\begin{grammarlens}[title={What you've built}]
One more word for everyday life at A2.
\end{grammarlens}

\subsection*{Your turn}
\begin{itemize}
\raggedright
  \item \emph{Say it:} \emph{camakvũ}
  \item \emph{Say it:} \emph{camakvũ}, once more
  \item \emph{Say it:} \emph{ākāśmā̃ vījḷī camakśe} — Lightning will flash in the sky.
  \item \emph{Recall:} say the Gujarati for daylight, then the Gujarati for to thunder, then say \emph{camakvũ} again
\end{itemize}

\begin{tcolorbox}[breakable,colback=teal!4,colframe=teal!35!black,title={Before you move on}]
What does \gu{ચમકવું} mean? (\textquotedblleft{}To shine, to flash\textquotedblright{}.) Say it once more.
\end{tcolorbox}
\section[\texorpdfstring{\gu{છવાવું} (chavāvũ)}{છવાવું (chavāvũ)}]{\gu{છવાવું} (chavāvũ) — to cloud over, to spread over}

\label{lesson:GU-C219-chavavu}



\begin{quote}
Before the new one: say the Gujarati for to thunder, then the Gujarati for to shine.
\end{quote}

\subsection*{You'll want to know: \gu{છવાવું}}
\textbf{\gu{છવાવું}} — \emph{chavāvũ} — \textquotedblleft{}to cloud over, to spread over\textquotedblright{}.

\begin{quote}
\textbf{\gu{સાંજે} \gu{આકાશમાં} \gu{વાદળ} \gu{છવાશે}.} — \emph{sā̃je ākāśmā̃ vādaḷ chavāśe} — \textquotedblleft{}Clouds will cover the sky in the evening.\textquotedblright{}
\end{quote}

\begin{cousinweb}
\textbf{\gu{છવાવું}} (\emph{chavāvũ}) is built on \emph{chāvũ}, \textquotedblleft{}to cover, to thatch\textquotedblright{}, with the \emph{-ā-} of a thing that happens by itself. CDIAL lists \emph{chāvũ}, \textquotedblleft{}to thatch\textquotedblright{}, under Sanskrit \emph{chādayati}, \textquotedblleft{}covers, veils\textquotedblright{} (CDIAL 5018).
\end{cousinweb}

\begin{grammarlens}[title={What you've built}]
One more word for everyday life at A2.
\end{grammarlens}

\subsection*{Your turn}
\begin{itemize}
\raggedright
  \item \emph{Say it:} \emph{chavāvũ}
  \item \emph{Say it:} \emph{chavāvũ}, once more
  \item \emph{Say it:} \emph{sā̃je ākāśmā̃ vādaḷ chavāśe} — Clouds will cover the sky in the evening.
  \item \emph{Recall:} say the Gujarati for to thunder, then the Gujarati for to shine, then say \emph{chavāvũ} again
\end{itemize}

\begin{tcolorbox}[breakable,colback=teal!4,colframe=teal!35!black,title={Before you move on}]
What does \gu{છવાવું} mean? (\textquotedblleft{}To cloud over, to spread over\textquotedblright{}.) Say it once more.
\end{tcolorbox}
\section[\texorpdfstring{\gu{કરા} (karā)}{કરા (karā)}]{\gu{કરા} (karā) — hailstones, hail}

\label{lesson:GU-C219-kara}



\begin{quote}
Before the new one: say the Gujarati for to shine, then the Gujarati for to cloud over.
\end{quote}

\subsection*{You'll want to know: \gu{કરા}}
\textbf{\gu{કરા}} — \emph{karā} — \textquotedblleft{}hailstones, hail\textquotedblright{}.

\begin{quote}
\textbf{\gu{કાલે} \gu{કરા} \gu{પડશે}.} — \emph{kāle karā paḍśe} — \textquotedblleft{}Hail will fall tomorrow.\textquotedblright{}
\end{quote}

Masculine, and usually plural.

\begin{cousinweb}
\textbf{\gu{કરા}} (\emph{karā}) is a tadbhava, a word inherited from Sanskrit and reshaped by sound change. It comes from \emph{karaka}, \textquotedblleft{}hail\textquotedblright{} (CDIAL 2782).
\end{cousinweb}

\begin{grammarlens}[title={What you've built}]
One more word for everyday life at A2.
\end{grammarlens}

\subsection*{Your turn}
\begin{itemize}
\raggedright
  \item \emph{Say it:} \emph{karā}
  \item \emph{Say it:} \emph{karā}, once more
  \item \emph{Say it:} \emph{kāle karā paḍśe} — Hail will fall tomorrow.
  \item \emph{Recall:} say the Gujarati for to shine, then the Gujarati for to cloud over, then say \emph{karā} again
\end{itemize}

\begin{tcolorbox}[breakable,colback=teal!4,colframe=teal!35!black,title={Before you move on}]
What does \gu{કરા} mean? (\textquotedblleft{}Hailstones, hail\textquotedblright{}.) Say it once more.
\end{tcolorbox}
\section[\texorpdfstring{\gu{માવઠું} (māvṭhũ)}{માવઠું (māvṭhũ)}]{\gu{માવઠું} (māvṭhũ) — unseasonal rain}

\label{lesson:GU-C219-mavthu}



\begin{quote}
Before the new one: say the Gujarati for to cloud over, then the Gujarati for hailstones.
\end{quote}

\subsection*{You'll want to know: \gu{માવઠું}}
\textbf{\gu{માવઠું}} — \emph{māvṭhũ} — \textquotedblleft{}unseasonal rain\textquotedblright{}.

\begin{quote}
\textbf{\gu{આ} \gu{વરસે} \gu{શિયાળામાં} \gu{માવઠું} \gu{થશે}.} — \emph{ā varase śiyāḷāmā̃ māvṭhũ thaśe} — \textquotedblleft{}This year there will be rain out of season in winter.\textquotedblright{}
\end{quote}

A neuter noun.

\begin{cousinweb}
\textbf{\gu{માવઠું}} (\emph{māvṭhũ}) is a tadbhava. It comes from \emph{*māghavṛṣṭi}, \textquotedblleft{}rain in the month of Māgha\textquotedblright{}, a reconstructed form built on Sanskrit \emph{vṛṣṭi}, \textquotedblleft{}rain\textquotedblright{}; CDIAL glosses the Gujarati word \textquotedblleft{}shower of rain at an unseasonable time\textquotedblright{} (CDIAL 9994).
\end{cousinweb}

\begin{grammarlens}[title={What you've built}]
One more word for everyday life at A2.
\end{grammarlens}

\subsection*{Your turn}
\begin{itemize}
\raggedright
  \item \emph{Say it:} \emph{māvṭhũ}
  \item \emph{Say it:} \emph{māvṭhũ}, once more
  \item \emph{Say it:} \emph{ā varase śiyāḷāmā̃ māvṭhũ thaśe} — This year there will be rain out of season in winter.
  \item \emph{Recall:} say the Gujarati for to cloud over, then the Gujarati for hailstones, then say \emph{māvṭhũ} again
\end{itemize}

\begin{tcolorbox}[breakable,colback=teal!4,colframe=teal!35!black,title={Before you move on}]
What does \gu{માવઠું} mean? (\textquotedblleft{}Unseasonal rain\textquotedblright{}.) Say it once more.
\end{tcolorbox}
